%%%
% src::
%     url = https://tex.stackexchange.com/a/634014/6880
%%%

\documentclass[border = 3pt]{standalone}

%%%
% Le \pack ''nicematrix'' est le couteau suisse des matrices.
%%%
\usepackage{nicematrix}

%%%
% Le \pack \tikz va être utilisé pour décorer les matrices.
%%%
\usepackage{tikz}

\begin{document}

%%%
% ''nicematrix'' fournit ''\hdashedline'' pour dessiner des
% \delims hachurées horizontales.
% Pour les hachures verticales, nous créons un nouveau style
% de ligne utilisable directement via un simple spécificateur.
%%%
\NiceMatrixOptions{
  custom-line = {
    letter      = I ,
    tikz        = dashed ,
    total-width = \pgflinewidth
  }
}

%%%
% L'emploi de ''first-row'' et ''first-col'' permet d'indiquer
% des étiquettes dans la \1ere ligne et la \1ere colonne du code
% \latex, ces étiquettes ne faisant pas partie de la matrice
% finale.
%%%
$H
=
 \begin{pNiceArray}{ccIcc}[
   first-row,
   first-col
 ]
     & a & b & c & d  \\
   e & 1 & 2 & 3 & 4  \\
   f & 5 & 6 & 7 & 8  \\
 \hdashedline
   g & 9 & 0 & 1 & 2  \\
   h & 3 & 4 & 5 & 7
 \end{pNiceArray}
$

\end{document}
